\chapter*{Preface}
\addcontentsline{toc}{chapter}{Preface}

Book 7 asked what a metric looks like once it is built to encode preference; the present book asks what time looks like once it is built to encode entropy. The two questions turn out to have a single answer. In the Relativistic Scalar-Vector Plenum, the passage of time is nothing more than the smoothing of curvature under entropy flow, and Ricci flow, the deformation of a manifold's metric by its own curvature, is here reinterpreted as the universal law of thermodynamic equilibration rather than as a purely mathematical technique for studying manifolds in the abstract. This book completes the series' transition from algebraic potential, developed across Cycle I, to dynamic geometry: the arrow of time itself emerges here as nothing more exotic than the entropy field's own irreversible diffusion across the manifold Book 6 constructed.

\chapter*{Diagrammatic Overview}
\addcontentsline{toc}{chapter}{Diagrammatic Overview}

Three constructions already introduced in this series converge in the present book: the entropic Ricci flow sketched in Book 6's Part III, the teleodynamic curvature of Book 7, and the BV antifield regularization of Book 5. This book does not introduce these as independent topics but shows how the flow that smooths curvature, the metric that encodes preference, and the antifield structure that keeps smoothing consistent are three aspects of a single dynamical process, here given its full development for the first time.

\chapter*{Reading Note}
\addcontentsline{toc}{chapter}{Reading Note}

Treat every evolution equation in this book as a story of reconciliation: each describes some quantity, curvature, a functional, a singularity, being brought back into alignment with the plenum's entropy budget, and the mathematics throughout is written with that reconciliation, rather than mere formal deformation, as its actual subject.

\part{Foundations of Ricci Flow in the Plenum}

\chapter{The Geometric Heat Equation}

Book 6's Chapter 5 already established the entropic Ricci flow $\partial_t g_{ij} = -2 R_{ij}(\SEntropy)$ as this series' replacement for Hamilton's ordinary Ricci flow $\partial_t g_{ij} = -2 R_{ij}$, and the present chapter does not re-derive that replacement; it takes the equation as given and develops what that earlier, foundational treatment had no occasion to supply, a diffusion equation for the entropy field itself driving the flow:
\[
\partial_t \SEntropy = \Delta \SEntropy + |\grad \SEntropy|^2.
\]
This equation identifies the flow parameter $t$ as the entropic parameter of smoothing directly: $t$ does not index an externally imposed time coordinate but is defined by, and only by, the rate at which this diffusion equation carries $\SEntropy$ toward equilibrium.

The analogy between heat diffusion and moral equilibration proposed here treats the diffusion equation's two terms as distinct modes of reconciliation: the Laplacian term $\Delta \SEntropy$ redistributes entropy locally, smoothing sharp gradients into gentler ones, while the quadratic term $|\grad \SEntropy|^2$ amplifies the diffusion wherever gradients are already steep, giving the equation a self-reinforcing character in regions of high tension that has a direct moral reading: reconciliation, on this account, accelerates precisely where it is most needed, not uniformly across the manifold.

